%%%%%%%%%%%%%%%%%%%%%%%%%%%%%%%%%%%%%%%%%%%%%%%%%%%%%%%%%%%%%%%%%
% Dissertacao de Mestrado / Dept Fisica, CFM, UFSC              %
% Andre@UFSC - 2011                                             %
%%%%%%%%%%%%%%%%%%%%%%%%%%%%%%%%%%%%%%%%%%%%%%%%%%%%%%%%%%%%%%%%%

%:::::::::::::::::::::::::::::::::::::::::::::::::::::::::::::::%
%                                                               %
%                          Cap�tulo 2                           %
%                                                               %
%:::::::::::::::::::::::::::::::::::::::::::::::::::::::::::::::%

%***************************************************************%
%                                                               %
%                            Galex                              %
%                                                               %
%***************************************************************%

\chapter{GALEX}
\label{sec:Galex}

O {\em Galaxy Evolution Explorer} (\galex) � um telesc�pio espacial de pequeno
porte da NASA\footnote{{\em NASA Small Explorer} ({\em SMEX}) -
http://explorers.gsfc.nasa.gov/missions.html}, lan�ado em 28 de abril de 2003
para conduzir um {\em survey} de todo o c�u numa faixa espectral do ultravioleta
(1350--2750\AA). O objetivo principal do \galex � estudar a evolu��o da taxa de
forma��o estelar em gal�xias \citep{Martin2005}. Os dados coletados pela miss�o
s�o publicados em {\em Data Releases} peri�dicos. O {\em Data Release} utilizado
neste trabalho � o {\em GR6 (GALEX Release 6)}.

 A miss�o consiste em uma s�rie de {\em surveys} fotom�tricos e
 espectrosc�picos (ver tabela \ref{tab:GalexSurveys}). As observa��es s�o feitas
 em duas bandas espectrais, ultravioleta distante ({\em far ultraviolet}, FUV),
 de 1350 a 1750\AA, e ultravioleta pr�ximo ({\em near ultraviolet}, NUV), de
 1750 a 2750\AA. A curva de transmiss�o dos filtros utilizados nessas bandas
 pode ser visto na figura \ref{fig:GalexFilters}. A espectroscopia � feita
 inserindo-se no caminho �tico um {\em grism}, que consiste num prisma combinado
 com uma rede de difra��o. Obt�m-se deste modo um espectro
 de baixa resolu��o para cada objeto no campo, conforme
 desctrito por \cite{Morrissey2007}.

\begin{table}
	\caption[{\em Surveys} realizados pelo \galex.]{Informa��es retiradas de
	\cite{Martin2005}.}
	\begin{tabular}{l r r}
		{\em Survey} & Cobertura do c�u ($graus^2$) & Magnitude AB limite \\ 
		\midrule
		{\em All-sky Imaging Survey (AIS)} & 26000 & 20.5 \\
		{\em Medium Imaging Survey (MIS)} & 1000 & 23 \\
		{\em Deep Imaging Survey (DIS)} & 80 & 25 \\
		{\em Nearby Galaxy Survey (NGS)} & 80 & 27.5 $arcsec^{-2}$  \\
		{\em Wide Field Spectroscopic Survey (WSS)} & 80 & 20 \\
		{\em Medium-deep Spectroscopic Survey (MSS)} & 8 & 21.5--23 \\
		{\em Deep Spectroscopic Survey (DSS)} & 2 & 23--24 \\
	\end{tabular}
	\label{tab:GalexSurveys}
\end{table}


% FIXME: Adicionar figura.
\begin{figure}
	\includegraphics[width=0.8\textwidth]{figuras/test.eps}
	\caption[Curvas de transmiss�o dos filtros do \galex.]
	{Curvas de transmiss�o dos filtros do \galex, medidas em
	laborat�rio \citep{Morrissey2005}.}
	\label{fig:GalexFilters}
\end{figure}

%***************************************************************%
%                                                               %
%                     Galex - Objetivos                         %
%                                                               %
%***************************************************************%

\section{Objetivos}
\label{sec:Galex:Objetivos}


Lalala \citep{Morrissey2007}


%***************************************************************%
%                                                               %
%                Galex - O c�u no ultravioleta                  %
%                                                               %
%***************************************************************%

\section{O c�u no ultravioleta}
\label{sec:Galex:CeuUV}

%***************************************************************%
%                                                               %
%                     Galex - Resultados obtidos                %
%                                                               %
%***************************************************************%

\section{Resultados obtidos}
\label{sec:Galex:Resultados}

Galex papers.

%***************************************************************%
%                                                               %
%                     Galex - Banco de dados                    %
%                                                               %
%***************************************************************%

\section{Banco de dados}
\label{sec:Galex:BancoDeDados}



% End of this chapter
